% Style based on the supplied Waseem_Generic_AI_ML_GenAI_AgenticAI.tex.
% Python replaces template tokens. Edit layout here, content in JSON/portfolio JS.
\documentclass[9pt,letterpaper]{extarticle}
\usepackage[left=0.68in,right=0.68in,top=0.45in,bottom=0.45in]{geometry}
\usepackage[T1]{fontenc}
\input{glyphtounicode}
\pdfgentounicode=1
\usepackage[utf8]{inputenc}
\usepackage{lmodern}
\usepackage{microtype}
\usepackage{pifont}
\usepackage{xcolor}
\usepackage{enumitem}
\usepackage{titlesec}
\usepackage[hidelinks]{hyperref}
\usepackage{tabularx,array,parskip,ragged2e,needspace}
\pagestyle{empty}
\setlength{\parindent}{0pt}
\setlength{\parskip}{1pt}
\definecolor{accent}{HTML}{155B8A}
\titleformat{\section}{\large\bfseries\color{accent}}{}{0pt}{}[\vspace{0.5pt}\color{accent}\titlerule]
\titlespacing*{\section}{0pt}{3pt}{2pt}
% Compact subsections beneath the main accent-rule sections.
\titleformat{\subsection}{\normalsize\bfseries\color{accent}}{}{0pt}{}
\titlespacing*{\subsection}{0pt}{3pt}{1pt}
\setlist[itemize]{leftmargin=1.15em,itemsep=0.8pt,topsep=1pt,parsep=0pt,partopsep=0pt}
\newcommand{\role}[4]{%
  \Needspace{4\baselineskip}\textbf{#1}\hfill\textbf{#2}\\[-1pt]
  \textit{#3}\hfill\textit{#4}\par}
\newcommand{\entry}[2]{\Needspace{4\baselineskip}\textbf{#1}\hfill\textit{#2}\par}
\newcommand{\skillrow}[2]{\textbf{#1:} & #2\\}
\hypersetup{pdftitle={Waseem Raza, PhD - Healthcare AI \& Wearable Intelligence},pdfauthor={Waseem Raza, PhD}}
\begin{document}

\begin{center}
{\LARGE\bfseries Waseem Raza, PhD}{\normalsize\mdseries }\\[2pt]
{\large\bfseries Machine Learning Engineer \textbar{} Healthcare AI \& Wearable Intelligence}\\[1pt]
{\small Healthcare AI \textbar{} Wearable Sensing \textbar{} Deep Learning \textbar{} Reliable ML}\\[3pt]
{\footnotesize \ding{41}\enspace \href{mailto:waseem.placements@gmail.com}{waseem.placements@gmail.com} \textbar{} \href{mailto:waseemraza844@gmail.com}{waseemraza844@gmail.com}\\[1pt]New Jersey, United States \quad \ding{37}\enspace \href{tel:+17328039703}{+1 732-803-9703}\\[1pt]\href{https://www.linkedin.com/in/waseem-raza-phd-9a737858/}{LinkedIn} \textbar{} \href{https://github.com/WaseemRaza844}{GitHub} \textbar{} \href{https://scholar.google.com/citations?user=d58Drr0AAAAJ\&hl=en}{Google Scholar} \textbar{} \href{https://waseemraza844.github.io/portfolio/learning.html}{Coursera Courses} \textbar{} \href{https://waseemraza844.github.io/portfolio/}{Portfolio}}
\end{center}
\vspace{-5pt}
\Needspace{3\baselineskip}
\section{Professional Profile}
Machine Learning Engineer and healthcare-AI researcher with a PhD in Electrical and Computer Engineering and experience applying deep learning, time-series modeling, wearable sensing, and reliable ML to real-world health-monitoring problems. Expertise includes Python, PyTorch, TensorFlow, scikit-learn, CNNs, LSTMs, attention, activity recognition, smartwatch IMU analytics, medication-taking detection, personalization, domain adaptation, few-shot learning, cross-user validation, and model evaluation. Complementary industry experience includes production ML, cloud data pipelines, data-quality validation, and multimodal analytics.
\Needspace{2\baselineskip}
\subsection{Technical Competencies}
\Needspace{2\baselineskip}\noindent{\small\textbf{Healthcare AI \& Wearables:} Wearable sensing, medication-adherence monitoring, activity recognition, smartwatch IMU data, human-centered sensing, and health-data applications}\par
\Needspace{2\baselineskip}\noindent{\small\textbf{Deep Learning:} CNNs, LSTMs, attention, transfer learning, domain adaptation, representation learning, classification, and sequence modeling}\par
\Needspace{2\baselineskip}\noindent{\small\textbf{Time-Series \& Sensor ML:} Windowing, segmentation, feature engineering, temporal modeling, cross-user evaluation, personalization, and continuous-event detection}\par
\Needspace{2\baselineskip}\noindent{\small\textbf{Reliable \& Personalized ML:} Distribution shift, few-shot adaptation, fairness-aware evaluation, robustness analysis, LOSO validation, and error analysis}\par
\Needspace{2\baselineskip}\noindent{\small\textbf{Data \& Engineering:} Python, PyTorch, TensorFlow, scikit-learn, Pandas, NumPy, cloud analytics, data pipelines, APIs, and reproducible experimentation}\par
\Needspace{2\baselineskip}\noindent{\small\textbf{Research \& Evaluation:} Experimental design, statistical analysis, sensitivity/specificity, precision/recall/F1, cross-validation, visualization, and technical communication}\par
\Needspace{4\baselineskip}\textbf{Leadership and Professional Strengths}\par
\begin{itemize}
\item \textbf{Ownership and accountability:} Lead AI/ML initiatives through requirements analysis, data preparation, model development, validation, stakeholder review, and deployment planning.
\item \textbf{Technical leadership:} Define technical approaches, evaluate alternatives, coordinate experimentation, and communicate findings to guide projects from ambiguous requirements to measurable outcomes.
\item \textbf{Cross-functional collaboration:} Work with software engineers, network specialists, researchers, geologists, and public-safety stakeholders to translate domain requirements into practical data and AI solutions.
\item \textbf{Continuous learning:} Apply ongoing study in generative AI, agentic AI, RAG, LLM evaluation, and cloud-native platforms through hands-on projects and professional coursework.
\item \textbf{Reliability and commitment:} Maintain reproducible workflows, careful validation, and attention to detail throughout experimentation and delivery.
\item \textbf{Adaptability and problem solving:} Combine critical thinking, systematic experimentation, and data-driven analysis to address unfamiliar technologies and complex operational problems.
\end{itemize}
\Needspace{3\baselineskip}
\section{Professional Experience}
\role{Machine Learning Engineer}{02/2025 - Present}{AT\&T, contractor via Innovatix Inc.}{Bedminster, NJ}
\begin{itemize}\interlinepenalty=10000
\item Lead end-to-end development and validation of production machine-learning workflows for high-impact operational decision support, including data preparation, model evaluation, error analysis, and reliability testing.
\item Design domain-adaptive neural models and reusable preprocessing pipelines to improve model generalization under heterogeneous data distributions.
\item Develop multimodal AI workflows that combine structured and unstructured operational data, providing transferable experience for complex real-world data environments.
\item Partner with cross-functional stakeholders to translate domain requirements into validated ML workflows and production-oriented data assets.
\end{itemize}
\role{Data Scientist and AI/ML Engineer}{09/2024 - 12/2024}{Impac Exploration Services}{Weatherford, OK}
\begin{itemize}\interlinepenalty=10000
\item Developed scalable analytical and ETL pipelines using Python, Azure Databricks, Azure Functions, and Synapse Analytics for large domain datasets.
\item Built reusable ingestion, cleansing, transformation, and validation workflows for machine-learning and LLM-assisted analytics.
\item Implemented cross-validation, data-quality, and consistency checks to improve reliability and reproducibility of analytical outputs.
\end{itemize}
\role{Advanced-Degrees Intern}{06/2024 - 08/2024}{AT\&T Labs}{Austin, TX}
\begin{itemize}\interlinepenalty=10000
\item Analyzed 500K+ LTE/5G-NR diagnostic logs and applied ML methods to detect anomalies, reducing diagnostic time by 20\%.
\item Investigated 5G NR-NTN cell-selection behavior and improved decision algorithms by 15\% through data-driven analysis and validation.
\item Conducted structured analysis of large operational datasets and communicated technical findings to engineering teams supporting next-generation wireless systems.
\end{itemize}
\role{Advanced Wireless Technology Development Intern}{09/2022 - 01/2023}{Futurewei Technologies}{Schaumburg, IL}
\begin{itemize}\interlinepenalty=10000
\item Improved autoencoder-based channel-state-information compression by 25\% using transfer learning and systematic experimentation on large wireless datasets.
\item Conducted 50+ experiments to validate model performance, robustness, and generalization under changing data conditions.
\item Collaborated with research and engineering stakeholders to translate modeling results into recommendations for next-generation communication systems.
\end{itemize}
\role{Graduate Research Assistant}{09/2019 - 07/2025}{AI4Networks, University of Oklahoma}{Tulsa, OK}
\begin{itemize}\interlinepenalty=10000
\item Designed reusable Python simulation and data-generation platforms for large-scale experimentation, including a 3GPP-compliant network simulator and an AI-enabled 5G testbed.
\item Built transfer-learning, adversarial-learning, GAN, autoencoder, and reinforcement-learning pipelines for sparse-data and cross-domain modeling.
\item Developed robustness evaluation frameworks for models operating under distribution shifts, missing data, and limited labeled samples; led 100+ experiments from problem formulation through validation.
\item Developed scalable self-healing and outage-management frameworks that improved modeled network resilience by 30\%.
\item Published peer-reviewed research, presented technical results internationally, and mentored junior researchers in ML implementation, experimental design, and technical writing.
\end{itemize}
\Needspace{3\baselineskip}
\section{Education}
\entry{PhD, Electrical and Computer Engineering}{09/2019 - 07/2025}
University of Oklahoma, Norman, Oklahoma, USA\par
\entry{MSc, Telecommunication Engineering}{08/2014 - 08/2016}
University of Engineering and Technology, Taxila, Pakistan\par
\entry{BSc, Telecommunication Engineering}{11/2010 - 06/2014}
University of Engineering and Technology, Taxila, Pakistan\par
\Needspace{3\baselineskip}
\section{Selected Projects \& Research}
\Needspace{2\baselineskip}
\subsection{Projects \& Applied Research}
\entry{Wearable Intelligence for Health Monitoring}{Research Project \textbar{} 2024-2026}
Developed deep-learning workflows for recognizing medication-taking and dietary activities from smartwatch accelerometer and gyroscope signals in free-living settings. Evaluated windowing, feature extraction, CNN/LSTM-style sequence models, cross-user generalization, and recall-oriented performance for health-monitoring applications.\par
{\small\textit{Tools \& Methods:} Wearable AI, Smartwatch IMU, Time Series, Deep Learning, CNN/LSTM, Feature Engineering, Cross-User Evaluation, Activity Recognition, Healthcare AI\par}
\vspace{1pt}
\entry{Cloud Computer Vision (GCP AutoML / Object Detection)}{Guided Project \textbar{} 2025}
Implemented cloud-based computer-vision workflows for image classification, object detection, and vehicle counting using managed model services and Python inference pipelines. Evaluated preprocessing, inference outputs, class/count aggregation, and model performance across visual-data tasks.\par
{\small\textit{Tools \& Methods:} Python, Computer Vision, Object Detection, Image Classification, Google Cloud AutoML, Streaming, Preprocessing, Inference, Model Evaluation\par}
\vspace{1pt}
\entry{OCR \& Document Intelligence Extraction Pipeline}{Guided Project \textbar{} 2025}
Implemented a document-intelligence pipeline that converts reports, invoices, and contract-style documents into structured outputs through OCR, layout/table parsing, entity and field extraction, normalization, and validation. Designed the workflow for downstream analytics and data integration rather than isolated text extraction.\par
{\small\textit{Tools \& Methods:} OCR, Document Intelligence, LLMs, Named Entity Recognition, Table Extraction, Structured JSON, Normalization, Validation, Data Integration\par}
\vspace{1pt}
\entry{\href{https://github.com/WaseemRaza844/agentic-ai-projects-by-was}{Agentic AI and RAG Applications}}{Portfolio Project \textbar{} 2026}
Built grounded GenAI workflows that combine document retrieval, vector search, stateful orchestration, tool-using agents, memory, evaluation, and multi-agent handoffs. Emphasized controllable execution, traceable context flow, and reliable responses for multi-step tasks rather than single unconstrained model calls.\par
{\small\textit{Tools \& Methods:} Python, LangGraph, LangChain, CrewAI, AutoGen, RAG, Vector Databases, Embeddings, Tool Calling, Agent Memory, Evaluation\par}
\vspace{1pt}
\entry{Amazon Bedrock Generative AI \& Agent Workflows}{Certification Project \textbar{} 2026}
Implemented foundation-model prompting, knowledge-base retrieval, RAG integration, model-customization exercises, workflow automation, and intelligent-agent patterns using Amazon Bedrock. Explored production-oriented cloud GenAI patterns with managed model access, retrieval-backed grounding, and reusable application workflows.\par
{\small\textit{Tools \& Methods:} Amazon Bedrock, Foundation Models, RAG, Knowledge Bases, LangChain, AI Agents, Prompt Engineering, AWS GenAI Services, Workflow Automation\par}
\vspace{1pt}
\entry{AI-Ready Lakehouse Pipeline}{Portfolio Project \textbar{} 2024}
Designed an AI-ready lakehouse pipeline that prepares structured and semi-structured data for analytics and downstream machine-learning workloads. Implemented repeatable ingestion, cleansing, transformation, validation, and scalable processing patterns suitable for reproducible cloud data workflows.\par
{\small\textit{Tools \& Methods:} Azure Databricks, Apache Spark, PySpark, ETL/ELT, Data Ingestion, Schema Transformation, Data Validation, Feature Preparation, Cloud Data Engineering\par}
\vspace{1pt}
\Needspace{2\baselineskip}
\subsection{Publications}
\entry{Real-World Medication Adherence Monitoring Using Hybrid Neural Networks and Smartwatches}{2026}
\textit{HealthINF \textbar{} Accepted}\par
\entry{\href{https://doi.org/10.3390/s25247651}{On Multi-Parameter Optimization and Proactive Reliability in 5G and Beyond Cellular Networks}}{2025}
\textit{Sensors 2025}\par
\entry{\href{https://doi.org/10.1109/PIMRC59610.2024.10817403}{An AI-Driven Framework for Enhancing Resilience in Propagation Models to Enable Digital Twin}}{2024}
\textit{IEEE PIMRC}\par
\entry{Data-Driven Intelligent Outage Management for High Shadowing Environments in 5G\&B Networks}{2024}
\textit{IEEE SmartNets}\par
\Needspace{3\baselineskip}
\section{Certifications \& Professional Development}
\Needspace{3\baselineskip}\noindent\textbf{1. Health Data Science Foundation} {\small --- \textcolor{blue}{\underline{\href{https://www.coursera.org/account/accomplishments/verify/42IDM121SWBS}{University of Illinois Urbana-Champaign / Coursera}}} \textbar{} Completed \textbar{} March 15, 2026}\par
\textbf{Focus:} Builds foundational skills for working with healthcare data by connecting data-science methods with clinical and health-domain questions. Emphasizes data representation, preprocessing, analysis workflows, interpretation, and practical constraints associated with health datasets.\par
{\small\textbf{Tools \& Skills:} health data science, healthcare datasets, data preprocessing, statistical analysis, data interpretation, reproducible analytics, health-domain reasoning\par}
\vspace{2pt}
\Needspace{3\baselineskip}\noindent\textbf{2. Generative AI for Healthcare} {\small --- \textcolor{blue}{\underline{\href{https://www.coursera.org/account/accomplishments/verify/P5AXL262N71P}{Google Cloud / Coursera}}} \textbar{} Completed \textbar{} February 17, 2026 \textbar{} \textcolor{blue}{\underline{\href{https://waseemraza844.github.io/portfolio/certificates/Coursera P5AXL262N71P.pdf}{Certificate}}}}\par
\textbf{Focus:} Explores generative-AI applications in healthcare, including information workflows, domain-aware application design, and opportunities to augment health-data analysis and communication. Emphasizes responsible adoption, usefulness, and evaluation within healthcare constraints.\par
{\small\textbf{Tools \& Skills:} healthcare GenAI use cases, prompt and application design, health-data workflows, responsible AI, domain-specific evaluation, AI adoption considerations\par}
\vspace{2pt}
\Needspace{3\baselineskip}\noindent\textbf{3. Detecting COVID-19 with Chest X-Ray using PyTorch} {\small --- \textcolor{blue}{\underline{\href{https://www.coursera.org/account/accomplishments/verify/MI6TUKWF37N2}{Coursera}}} \textbar{} Completed \textbar{} January 27, 2026}\par
\textbf{Focus:} Applies deep learning to chest X-ray image classification using PyTorch, covering image preprocessing, model training, validation, and interpretation in a healthcare context. Demonstrates an end-to-end computer-vision workflow for medical-image analysis.\par
{\small\textbf{Tools \& Skills:} PyTorch, medical image classification, convolutional neural networks, image preprocessing, training and validation, model evaluation, healthcare AI\par}
\vspace{2pt}
\Needspace{3\baselineskip}\noindent\textbf{4. Introduction to Deep Learning \& Neural Networks with Keras} {\small --- \textcolor{blue}{\underline{\href{https://www.coursera.org/account/accomplishments/verify/Y6NP4GWLS4DZ}{IBM / Coursera}}} \textbar{} Completed \textbar{} August 12, 2026}\par
\textbf{Focus:} Introduces deep-learning and neural-network fundamentals using Keras, covering network architecture, training, validation, and model improvement. Connects concepts such as activation functions and optimization with practical neural-network implementation and evaluation.\par
{\small\textbf{Tools \& Skills:} Keras, deep learning, neural networks, model training, validation, activation functions, optimization, model evaluation\par}
\vspace{2pt}
\Needspace{3\baselineskip}\noindent\textbf{5. Machine Learning} {\small --- \textcolor{blue}{\underline{\href{https://www.coursera.org/account/accomplishments/verify/VU3Y5TBLVLZT}{Stanford Online / Coursera}}} \textbar{} Completed \textbar{} March 29, 2020}\par
\textbf{Focus:} Covers foundational machine-learning methods for supervised and unsupervised problems, including regression, classification, model evaluation, regularization, clustering, and practical model-selection considerations. Emphasizes translating data into predictive models while controlling generalization error.\par
{\small\textbf{Tools \& Skills:} supervised learning, regression, classification, regularization, clustering, model evaluation, bias-variance tradeoffs, machine-learning workflows\par}
\vspace{2pt}
\Needspace{3\baselineskip}\noindent\textbf{6. Introduction to Healthcare Management} {\small --- \textcolor{blue}{\underline{\href{https://www.coursera.org/account/accomplishments/verify/8CIQHG3YD4HM}{Coursera}}} \textbar{} Completed \textbar{} March 3, 2026}\par
\textbf{Focus:} Introduces management concepts for healthcare organizations, including organizational structure, stakeholders, service delivery, operations, leadership, and resource coordination. Emphasizes how managerial decisions affect quality, efficiency, and performance in healthcare systems.\par
{\small\textbf{Tools \& Skills:} healthcare management, healthcare operations, organizational structure, stakeholder management, resource coordination, service delivery, performance improvement\par}
\vspace{2pt}
\Needspace{3\baselineskip}\noindent\textbf{7. Python for Data Science, AI \& Development} {\small --- \textcolor{blue}{\underline{\href{https://www.coursera.org/account/accomplishments/verify/WRA89JEJU8EE}{IBM / Coursera}}} \textbar{} Completed \textbar{} July 12, 2026 \textbar{} \textcolor{blue}{\underline{\href{https://waseemraza844.github.io/portfolio/certificates/cloud-devops/C7. Coursera WRA89JEJU8EE.pdf}{Certificate}}}}\par
\textbf{Focus:} Builds practical Python skills for data science and AI development, including core programming, data structures, data manipulation, and application-oriented analytical workflows. Emphasizes using Python as a reusable foundation for data processing, modeling, and AI applications.\par
{\small\textbf{Tools \& Skills:} Python, data structures, data manipulation, Pandas concepts, NumPy concepts, data-analysis workflows, AI application development\par}
\vspace{2pt}
\end{document}
